\section{Executive summary}

The second assignment of this course has enhanced our previous implementation of the Catan simulator to accommodate **human gameplay**. New features of this enhanced version are the implementation of a **command-line interface** for human players to enter their commands, the implementation of the **robber** function, which is activated when the rolled numbers are 7, and the implementation of **JSON export** for visualization of the game. This version of the simulator has been enhanced to **pause** between human turns, waiting for the human to enter "Go" to continue the game. A total of **18 unit tests** were written before this implementation, ensuring the detection of bugs early in the process. This version of the simulator has been successfully enhanced to accommodate human gameplay while retaining all original functionality.

\section{Requirements traceability}

\begin{table}[h]
    {\small
    \begin{tabular}{@{}p{2cm}p{2cm}p{2cm}p{8cm}@{}}
        \toprule
        \multicolumn{1}{c}{\textbf{Req ID}} & \multicolumn{1}{c}{\textbf{Status}} & \multicolumn{1}{c}{\textbf{Implemented in}} & \multicolumn{1}{c}{\textbf{Design considerations}} \\
        R1.1 & \{Missing, Partial, Implemented\} & \code{MapSetup.java} & The software sets up the same map by a hard-wired specification. The specification is given in the \code{MapSetup} class' \code{tiles} data structure and used in the \code{generateMap()} method at the initialization phase of the program. \\
        ... & ... & ...  & ... \\ \bottomrule
    \end{tabular}}
\end{table}

\id{The goal of this matrix is to aid your TA verify the satisfaction of requirements. If the table is superficial or too verbose, your TA will have a hard time verifying the completeness of your deliverable, and that will result in lost marks.
Use the Mising-Partial-Implemented status markers in your final report. Use the code{} command to format file names and Java code elements, such as class names, method names, etc.}

\section{Getting ready for change}

\textbf{How many unit tests did you end up writing?}

\textbf{If you measured coverage, which tool did you use and what was you coverage?}

\section{Changing and extending your design}

\textbf{What were the main structural changes in your original design? Was it easy to localize the change and assess its impact compared to directly changing the code (without messing around with the design)?}

\textbf{Why was the automaton-based modelling useful and what benefits does it provide compared to the specification in the attached PDF? How did you translate the automaton-based model to program code?}

\section{Revise the program code and finish the functionality}

\textbf{Which parts of your system did break and how did you resolve those issues?}

\textbf{What was the benefit of having unit tests at this stage?}

\textbf{What were some of the challenges in these iterations (technical or conceptual)? What tools or integrated tool chains would help you in those steps? (You can specify a few abstract ideas, you don’t have to name specific tools or techniques.)}

\textbf{Elaborate on the OO mechanisms and SOLID principles in your implementation. Are these any different from the ones used in the design? If they are: why?}

\section{Roles and responsibilities}

\begin{itemize}
    \item {\textbf{Raadhikka Gupta} }
    \item {\textbf{Ranica Chawla} }
	\item {\textbf{Nitya Patel} }
	\item {\textbf{Krisha Patel} }
\end{itemize}
